\section{Computations}
\label{app:computations}

\emph{Certificates} check finitely many polynomial identities exactly,
over polynomial rings such as $\mathbb{F}_2[q]$ or $\mathbb{F}_2[q,\lambda]$;
they are part of the proofs of the statements that cite them, and
multilinearity extends them to all inputs. \emph{Experiments} compute
homological invariants in finitely many degrees, over finite fields or over
$\mathbb{F}_2(q,H_1,H_2)$; they are evidence with the stated scope and prove
nothing about other degrees or other fields. All scripts and their saved
outputs are in the directory \texttt{computations/} of the repository described
in \Cref{app:ai-declaration}. A rerun of the
scripts, recorded in the repository, executed 45 commands with Python~3.14 and
SageMath~10.10, and their mathematical output agreed with the saved results.
Three finite-field runs for $\Lambda_0$ were not repeated, and two stages for
$\Lambda_1$ and two stages of the attempt using minimal bimodule resolutions
over ${T\otimes_\kk T}$ were rerun from
saved intermediate results. The finite fields
$\mathbb{F}_{2^{16}}$ below are defined by the modulus
${x^{16}+x^5+x^3+x^2+1}$, and the parameters $q$, $H_1$ and $H_2$ are powers of a
generator, with exponents recorded in the saved outputs.

\subsection*{The cochain}

The following table, reproduced from the appendix of~\cite{OAI26ar}, defines the
cochain $p$ of \Cref{subsec:cocycle}. Lower-case letters denote the basis of $C$
of \Cref{subsec:algebra-C} and capital letters the dual basis of $DC$; an entry
$abcw$ in the row $q^d$ means that the coefficient of the basis vector $w$ in
$p(a,b,c)$ is $q^d$, and all other coefficients are zero. A script confirms that
the table below and the source table have the same $179$ entries
in the same order.

\begin{center}
\small
\begin{tabular}{r@{\quad}l}
 coefficient & input and output letters \\ \hline
 $q^2$ & \texttt{EuNT EyZX FNtU JtXZ NtEU NtXV XzJT ZxEY tEun tXuf} \\
 $q^2$ & \texttt{tXvn tXzj tuFn tuNf tvNn uFNT vNNT vtJu xEyz zJtx} \\
 $q$ & \texttt{EuUE EuVX EvVE ExXE EyYE EzZE FNnF FVvF FVxU FtTF} \\
 $q$ & \texttt{FtYU Ftun Ftxj JFtX JNtZ JUyX JVyZ JjEE JjTT JjXX} \\
 $q$ & \texttt{NFtU NNtV NUyU NVyV NjJN NjZV NnNN NnVV Ntuf Ntvn} \\
 $q$ & \texttt{Ntzj TtEE TtXX UuFF UuUU UvNF UvVU UxXU Uyun Uyxj} \\
 $q$ & \texttt{UzZU VvNN VvVV VxEU VxXV Vyuf Vyvn Vyzj VzZV XxEE} \\
 $q$ & \texttt{XxXX YuFT YuUY YvNT YvVY YxXY YyEE YyJJ YyTT YyXX} \\
 $q$ & \texttt{YyYY YyZZ YzZY ZxXZ ZzEE ZzJJ ZzTT ZzXX jJFF jJUU} \\
 $q$ & \texttt{jJjj jJnn jYyj jZvn jZzj nNFF nNUU ntTn tJjt tTNN} \\
 $q$ & \texttt{tTVV tTtt tYyt tZzt tzJn uFVX uFty uNFT uNUY uNte} \\
 $q$ & \texttt{uUyy uVye ujJu vNUE vNVX vNty vVyy vtTv xEuv xJFT} \\
 $q$ & \texttt{xJUY xJjx xJte xTtx xXuu xXvv xXyy xXzz xYyx xZye} \\
 $q$ & \texttt{xZzx xuFv xuNu xvNv yxJu yzJv zJNT zJVY zYyz zZuu} \\
 $q$ & \texttt{zZvv zZyy zZzz} \\
 $q^3$ & \texttt{EyJT JtEY tEyj tyJf} \\
 $1$ & \texttt{JNjX JVzX NNjU NVzU NjYU Njun Njxj Ntxt TNtX TVyX} \\
 $1$ & \texttt{Vyxt VzYU Vzun Vzxj YuNJ YuVZ jYun juNn jxJn txTn} \\
 $1$ & \texttt{uNNJ uNVZ uNjy uVzy xJNJ xJVZ xTNT xTVY yxTv zYuv} \\
 $1$ & \texttt{zuNv zxJv} \\
\end{tabular}
\end{center}

\subsection*{Certificates for \texorpdfstring{\Cref{sec:ar-route}}{Section 9}}

\begin{enumerate}
  \item \emph{The algebras $C$ and $T$}
  (\Cref{lemma:algebra-C,lemma:T-symmetric};
    \texttt{08-D-E/foundations\_}\allowbreak\texttt{certificate.py},
    over $\mathbb{F}_2[q]$). All
    $1000$ associators of basis vectors of $C$ and all $8000$ of $T$ vanish; the
    grading is respected by all products; the trace form of $T$ is symmetric,
    associative and has the stated dual bases; ${(ut)^2}={q(1+q)z}$. An earlier,
    independently written script (\texttt{02-ar-finite-data/}) checks the same
    identities over $\mathbb{F}_2(q)$ and $\mathbb{F}_{2^{16}}$ and computes the
    resolution of $s$ in degrees up to eight, in agreement with
    \Cref{prop:C-resolution}.
  \item \emph{The cochain} (\Cref{lemma:cocycle};
    \texttt{08-D-E/finite\_cochain\_certificate.py}, over $\mathbb{F}_2[q]$ with
    the coefficients of $\lambda^0$ and $\lambda^1$ compared separately, hence
    over $\mathbb{F}_2[q,\lambda]$). Every entry has radical inputs, matching
    corners and weight ${-1}$; the cocycle identity holds on all ${18^4=104\,976}$
    words of four radical basis vectors; the boundary identity holds on all
    $324$ pairs; ${p(t,x,j^*)=0}$, ${p(t,y,j^*)=q^3f}$, and the bar chain $\zeta$
    of \Cref{prop:cocycle-generator} is a cycle. The script
    \texttt{08-D-E/certificate\_receipt.py} confirms that its table equals the
    source table.
  \item \emph{Casimir identities} (\Cref{prop:lifts}; the same script). The sum
    \[
      \sum_wh_\lambda(w)w^*
    \]
    over the twenty basis vectors vanishes, and the sum
    over the radical basis vectors with left idempotent $r$ equals $r^*$, for
    ${r\in\set{e,f}}$. The other identities used in that proof are proved in the
    text.
\end{enumerate}

\subsection*{Dimensions}

The dimension of $C_1$ in \Cref{rem:ar-sizes} is computed exactly from the
corner dimensions of $T$ and $\operatorname{rad}T$, which give the dimensions
of the terms of the bar resolution, and from the cohomology $E$, $E^2$ and $E$
of $K$ in degrees $0$, $2$ and $4$ (\texttt{09-sizes/dimensions.py}). The
terms $K^0,\dots,K^4$ have dimensions $1\,761\,405\,952$, $148\,453\,376$,
$11\,713\,792$, $818\,688$ and $43\,264$, which gives ${\dim
C_1=1\,623\,889\,344}$
and
\[
  \begin{aligned}
    \dim_\kk\Lambda&=1600+159999^5\cdot1\,623\,889\,344\\
    &=170\,271\,818\,183\,326\,072\,615\,867\,045\,851\,312\,256.
  \end{aligned}
\]
The computation (\texttt{06-two-factor/}, over $\mathbb{F}_{2^{16}}$)
used minimal bimodule resolutions of the cones and determined, from the
dimensions of the terms and of the homology of the complex, cokernels of
dimensions $784\,704$ and $1\,377\,984$; it
stopped, before forming the bimodule matrices, at a linear system with
$3\,591\,168$ unknowns per right-hand side, beyond the bound of ${2\cdot10^6}$
set in advance. The same computation gives the dimensions of the Tate cohomology
groups in \Cref{prop:two-cone-profile} in degrees ${-4\leq a\leq7}$ for its sample
parameters.

\subsection*{Experiments for \texorpdfstring{\Cref{sec:obstructions}}{Section
10}}

\begin{enumerate}
  \item \emph{The algebra $\Lambda_0$} (\texttt{03-testbed-lambda0/}). Three
    runs over $\mathbb{F}_{2^{16}}$ in degrees up to six and one exact run over
    $\mathbb{F}_2(q,H_1,H_2)$ in degrees up to two compute minimal projective
    resolutions of $Z_0$ and the groups $\operatorname{Ext}^a(Z_0,Z_0)$ and
    $\operatorname{Ext}^a(Z_0,\Lambda_0)$. A separate script compares the
    surviving class with the cokernel of $\delta^0$; this map has rank one and
    kernel spanned by $(1,1)$.
  \item \emph{The algebra $\Lambda_1$ constructed over $T$}
  (\texttt{05-candidate1/}).
    Two runs over $\mathbb{F}_{2^{16}}$ construct minimal bimodule covers, the
    cocycle, the cone of dimension $312$, normalised lifts, the fibre of
    dimension $352$, the algebra $\Lambda_1$ and the module $Z_1$, and compute
    the dimensions of the Tate cohomology groups of the cone in degrees ${-4\leq
    a\leq4}$ (nonzero only in degrees
    $0$ and $1$, as in \Cref{prop:one-cone}), the matrix of $\delta^0$, and
    $\operatorname{Ext}^a$ for ${a=1,2,3}$.
  \item \emph{Further attempts using a single cone}
  (\texttt{07-one-factor-search/},
    over $\mathbb{F}_{2^{16}}$). For the six-dimensional trivial extension of
    ${\kk(1\to2)}$, its tensor product with $\kk[\epsilon]/(\epsilon^2)$, and the
    trivial extension of $T$, minimal resolutions of the simple modules give
    Ext dimensions up to degree twelve, directly for the first two and, for
    the third, directly in degrees $0$ and $1$ and by a tensor-product formula
    in degrees $2$ to $12$. None has polynomial Ext algebra on a generator of
    degree three. The same directory contains the six-dimensional example
    mentioned in \Cref{rem:weights}.
  \item \emph{Brackets for $T$} (\texttt{04-toda-bracket/}, over
    $\mathbb{F}_2(q)$ and over $\mathbb{F}_{2^n}$ for ${n}={8,12,16}$). An explicit
    defining system shows ${\langle\tau,\beta_0,\beta_0\rangle}={\set{0}}$, in
    agreement with \Cref{thm:tate-obstruction}.
\end{enumerate}

\subsection*{Checks of signs and finite cases}

The remaining scripts check signs and finite cases of general arguments; they
are not needed for the proofs. They cover the finite quotients of $G$ for
${m\leq10}$
(\texttt{D-C-selection-checks.py}); the signs of the bar decomposition, of the
simulation and two examples, including
\Cref{ex:derived-powers} (\texttt{D-A-checks.py}); the rectification for an
algebra with a nonzero null-homotopic relation
(\texttt{D-B-rectification-check.py});
and the signs of the cone in the proof of \Cref{thm:conversion}
(\texttt{D-D-cone-signs.py}).
